\documentclass[11pt,a4paper]{article}
\usepackage[utf8]{inputenc}
\usepackage[margin=1in]{geometry}
\usepackage{hyperref}
\usepackage{array}
\usepackage{booktabs}
\usepackage{enumitem}
\usepackage{titlesec}
\usepackage{xcolor}

\definecolor{cudenvergold}{HTML}{8A724F}
\definecolor{darkgrey}{HTML}{333333}

\titleformat{\section}{\large\bfseries\color{darkgrey}}{\thesection}{1em}{}[{\color{cudenvergold}\titlerule[1.5pt]}]
\titleformat{\subsection}{\normalsize\bfseries\color{darkgrey}}{\thesubsection}{1em}{}

\hypersetup{
    colorlinks=true,
    linkcolor=darkgrey,
    filecolor=magenta,      
    urlcolor=cudenvergold,
    pdftitle={ISMG 2050 Fall 2026 Syllabus},
    pdfpagemode=FullScreen,
}

\begin{document}

\begin{center}
    {\LARGE \bfseries University of Colorado Denver} \\
    \vspace{2mm}
    {\Large \bfseries Business School} \\
    \vspace{4mm}
    {\color{cudenvergold}\hrule height 2pt}
    \vspace{3mm}
    {\large \bfseries ISMG 2050: Business Problem Solving Tools} \\
    {\large \bfseries Course Syllabus --- Fall 2026} \\
\end{center}

\vspace{0.5cm}

\section*{Course Logistics \& Instructor Information}
\begin{itemize}[leftmargin=*]
    \item \textbf{Instructor:} Franklin Diaz, Adjunct Faculty
    \item \textbf{Email:} \href{mailto:franklin.diaz@ucdenver.edu}{franklin.diaz@ucdenver.edu} (Canvas Inbox preferred)
    \item \textbf{Class Schedule:} Monday \& Wednesday, 3:30 PM -- 4:45 PM Mountain Time (MT)
    \item \textbf{Classroom Location:} CU Denver Business School Building, Room 2501
    \item \textbf{Office Location:} Business School Building or Virtual via Zoom (by appointment)
    \item \textbf{Office Hours:} By appointment via Canvas Inbox
    \item \textbf{Student Success Coach (SSC):} Abdelkarim Mansour (\href{mailto:abdelkarim.mansour@ucdenver.edu}{abdelkarim.mansour@ucdenver.edu})
    \item \textbf{Term Dates:} Fall 2026 (08/24/2026 -- 12/12/2026)
    \item \textbf{Delivery Mode:} In-person meetings with online assignments and materials via Canvas
    \item \textbf{Prerequisites / Co-requisites:} None
\end{itemize}

\section*{Course Description}
ISMG 2050 introduces students to the technology and problem-solving skills necessary for success both at school and in the business world. This hands-on course focuses on making data-driven decisions using spreadsheets, data analysis, and visualization tools. Students solve quantitative and strategic problems across statistics, accounting, finance, marketing, management, and information systems. We specifically utilize Microsoft Excel and Tableau to equip students with practical methodologies for the modern business landscape.

As a business core course, a grade of \textbf{C-} or better (70.00\%) must be earned to satisfy Business graduation requirements and advance to higher-level business courses.

\section*{Course Goals \& Learning Objectives}
By the end of this course, you will be able to:
\begin{itemize}[noitemsep]
    \item Obtain foundational and intermediate proficiency in Microsoft Excel.
    \item Construct formulas utilizing mathematical operators, logical functions, and lookup arrays.
    \item Present data effectively with professional formatting, dynamic charts, and PivotTables.
    \item Apply logic models, What-If analysis, and financial functions to business decisions.
    \item Master ETL (Extract, Transform, Load) processes using Excel Power Query.
    \item Obtain proficiency in Tableau to build interactive dashboards and visual stories from multidimensional datasets.
\end{itemize}

\section*{Required Materials \& Technology}
\begin{itemize}[leftmargin=*]
    \item \textbf{Required Textbook:} \textbf{None.} There are no textbooks to purchase. All exercise guides, instructional notes, datasets, and walkthrough videos are provided at zero cost directly via Canvas modules.
    \item \textbf{Software \& Computer Lab Access:}
    \begin{itemize}[noitemsep]
        \item \textbf{Microsoft 365 (Excel):} Provided free via CU Denver student accounts for personal laptops.
        \item \textbf{Tableau Analytics:} Due to software compatibility and licensing constraints, \textbf{all Tableau coursework, instruction, and project work will be conducted directly using the Business School Computer Labs}, where full, pre-configured enterprise desktop installations are provided.
    \end{itemize}
    \item \textbf{Personal Computing Device:} A laptop (Windows or macOS) for routine classroom notes, Excel modeling, and Canvas submissions.
\end{itemize}

\section*{Grading Architecture}
Final grades are computed based on the following breakdown:
\begin{center}
    \begin{tabular}{lc}
        \toprule
        \textbf{Assessment Component} & \textbf{Weight} \\
        \midrule
        Hands-on Excel \& Tableau Assignments & 65\% \\
        Exams (Midterm and Final) & 20\% \\
        Tableau Capstone Project & 10\% \\
        Participation \& In-Class Engagement & 5\% \\
        \bottomrule
    \end{tabular}
\end{center}

\subsection*{Undergraduate Letter Grade Scale}
\begin{center}
    \begin{tabular}{llp{7.5cm}}
        \toprule
        \textbf{Letter Grade} & \textbf{GPA Value} & \textbf{Percentage Range} \\
        \midrule
        A  & (4.0) & 93.00\% -- 100.00\% (Superior/Excellent mastery) \\
        A- & (3.7) & 90.00\% -- 92.999\% \\
        B+ & (3.3) & 87.00\% -- 89.999\% \\
        B  & (3.0) & 84.00\% -- 86.999\% (Good/Better than average) \\
        B- & (2.7) & 80.00\% -- 83.999\% \\
        C+ & (2.3) & 77.00\% -- 79.999\% \\
        C  & (2.0) & 74.00\% -- 76.999\% (Competent/Average) \\
        C- & (1.7) & 70.00\% -- 73.999\% (Minimum Passing for Business Core) \\
        D+ & (1.3) & 67.00\% -- 69.999\% \\
        D  & (1.0) & 64.00\% -- 66.999\% \\
        D- & (0.7) & 60.00\% -- 63.999\% \\
        F  & (0.0) & 0.00\% -- 59.999\% (Failing) \\
        \bottomrule
    \end{tabular}
\end{center}
\textit{Note: CU Denver Business School guidelines target an undergraduate class GPA average between 2.3 (C+) and 3.0 (B). If necessary, grading thresholds may be adjusted to align with institutional standards.}

\subsection*{Late Submission Policy}
All assignments are due by \textbf{11:59 PM MT} on the designated due dates in Canvas. Late submissions incur a deduction of \textbf{10\% per calendar day late}. Assignments are not accepted after 9 days late (resulting in a score of \textbf{0} on day 10).

\subsection*{Mastery Retry Policy}
If you score below a C (74\%) on any assignment worth 5\% or more of the course grade, you may submit a mastery retry for partial credit. You must meet with the instructor or Student Success Coach, review the conceptual knowledge gaps, and resubmit the revised assignment within 7 days of receiving your original score (maximum score capped at 90\%).

\section*{AI Collaboration \& Academic Integrity Policy}
Artificial Intelligence tools (such as ChatGPT, Claude, Microsoft Copilot, and Gemini) are integral components of the modern analyst workflow. \textbf{Students are encouraged to actively use AI as a learning accelerator, debugging partner, and cognitive tutor.}

\subsection*{Encouraged Uses}
\begin{itemize}[leftmargin=*]
    \item \textbf{Function \& Formula Deconstruction:} Asking AI to break down complex nested formulas (e.g., INDEX/XMATCH or multi-criteria IF statements) into plain English.
    \item \textbf{Syntax \& Error Debugging:} Pasting broken formulas or error codes (\#VALUE!, \#REF!, \#N/A) to identify syntax mistakes or array mismatches.
    \item \textbf{Visualization Design:} Brainstorming chart types, color palettes, and dashboard layout ideas.
    \item \textbf{Concept Review:} Prompting AI to generate practice questions and conceptual drills to prepare for exams.
\end{itemize}

\subsection*{Prohibited Uses}
\begin{itemize}[leftmargin=*]
    \item Submitting raw AI-generated files, spreadsheets, or dashboards without completing and understanding the work yourself.
    \item Using AI or external resources during closed-book, in-class midterm and final examinations.
    \item Submitting unverified AI outputs or misrepresenting another person's work as your own.
\end{itemize}

\subsection*{AI Citation Requirement}
When AI is consulted to help solve an assignment problem, include an inline disclosure note in your submission:
\begin{quote}
    \textit{AI Disclosure: Used [Tool Name, e.g., ChatGPT / Gemini] to [Purpose, e.g., ``debug nested XLOOKUP formula and explain array height errors'']}.
\end{quote}
If no AI was used: \textit{AI Disclosure: None.}

\section*{Course Schedule \& Modular Breakdown}
\begin{itemize}[leftmargin=*]
    \item \textbf{Module 1: Creating a Worksheet and a Chart (Weeks 1 \& 2)}
    \begin{itemize}[noitemsep]
        \item \textit{Topics:} Ribbon navigation, grid mechanics, basic arithmetic functions, cell styles, proportional Pie Charts.
        \item \textit{Hands-on Projects:} Real Estate Budget, Dynamic Updates, Style Application, Laptop Analysis.
    \end{itemize}

    \item \textbf{Module 2: Formulas, Functions, and Formatting (Weeks 3 \& 4)}
    \begin{itemize}[noitemsep]
        \item \textit{Topics:} Flash Fill, math operators, conditional formatting rules, page setup.
        \item \textit{Hands-on Projects:} Salary Report, Cost Analysis, Customer Tracking, College Cost Calculator.
    \end{itemize}

    \item \textbf{Module 3: Large Worksheets, What-If Analysis, and Charting (Weeks 5 \& 6)}
    \begin{itemize}[noitemsep]
        \item \textit{Topics:} Absolute vs. Relative referencing (\$A\$1), nested IF logic, Sparklines, What-If analysis.
        \item \textit{Hands-on Projects:} Financial Projections, Logic Tests, Census Data.
        \item \textit{Assessment:} \textbf{Quiz 1 (Week 6).}
    \end{itemize}

    \item \textbf{Module 4: Financial Functions, Data Tables, and Amortization (Weeks 7 \& 8)}
    \begin{itemize}[noitemsep]
        \item \textit{Topics:} Time Value of Money (PMT, PV, FV), cash flow signs, amortization schedules, sheet protection.
        \item \textit{Hands-on Projects:} Mortgage Payment Calculator, Loan Analysis, Retirement Planning.
    \end{itemize}

    \item \textbf{Module 5: Multiple Worksheets and Advanced Function Library (Weeks 9 \& 10)}
    \begin{itemize}[noitemsep]
        \item \textit{Topics:} Multi-sheet 3-D formulas, data consolidation, custom formatting codes.
        \item \textit{Hands-on Projects:} Payroll Consolidation, Custom Codes, Multi-year Trends.
    \end{itemize}

    \item \textbf{Module 6: Creating, Sorting, and Querying Tables (Weeks 11 \& 12)}
    \begin{itemize}[noitemsep]
        \item \textit{Topics:} Structured table objects, XLOOKUP querying, Advanced AutoFilters.
        \item \textit{Hands-on Projects:} Account Management, Table Logic.
        \item \textit{Assessment:} \textbf{Midterm Examination (Week 12).}
    \end{itemize}

    \item \textbf{Module 7: Power Query (ETL), PivotTables, and Macros (Weeks 13 \& 14)}
    \begin{itemize}[noitemsep]
        \item \textit{Topics:} Power Query ETL data pipelines, PivotTables, PivotCharts, Macro automation.
        \item \textit{Hands-on Projects:} Power Query Transformation Pipeline, Interactive Slicers.
    \end{itemize}

    \item \textbf{Module 8: Tableau Visual Analytics and Dashboards in Computer Lab (Weeks 15 \& 16)}
    \begin{itemize}[noitemsep]
        \item \textit{Topics:} Business School Computer Lab sessions: Dimensions vs. Measures, interactive filters, dashboard storytelling.
        \item \textit{Deliverables:} \textbf{Tableau Capstone Project Due; Comprehensive Final Exam (Week 16).}
    \end{itemize}
\end{itemize}

\section*{Course Policies \& Institutional Disclosures}
\begin{itemize}[leftmargin=*]
    \item \textbf{Canvas Messaging:} Use the \textbf{Canvas Inbox} for all course correspondence. Email spam filters and firewalls frequently block macro-enabled (.xlsm) or packaged Tableau (.twbx) workbooks attached to direct emails.
    \item \textbf{Student Success Coach (SSC):} Abdelkarim Mansour (\href{mailto:abdelkarim.mansour@ucdenver.edu}{abdelkarim.mansour@ucdenver.edu}) is available for academic coaching. Meeting with the SSC to remediate an assignment with a grade of C or below restores \textbf{5\% of available points} on that assignment.
    \item \textbf{Classroom \& Lab Etiquette:} Laptops should be used for active in-class spreadsheet exercises; lab workstations are reserved for official course analytics tools. Mobile phone use and texting during class are prohibited.
    \item \textbf{Visible ID Badge:} CU Denver safety policy requires all students, faculty, and staff in the Business School Building to visibly display their university ID badge.
    \item \textbf{Accommodations (DRS):} Students requiring accommodations must register with Disability Resources and Services (Student Commons 2116, 303-315-3510, \href{mailto:disabilityresources@ucdenver.edu}{disabilityresources@ucdenver.edu}).
    \item \textbf{Counseling \& Support:} Free and confidential mental health support is available at the Counseling Center (Tivoli 454, 303-315-7270). Safety concerns may be reported to the CARE Team (303-315-7306).
\end{itemize}

\vfill
\noindent\rule{\textwidth}{0.4pt}
\small
\textit{Land Acknowledgement: We acknowledge the Cheyenne, Arapaho, and Ute nations as the original stewards of the land on which this university stands.}\\
\textit{This syllabus is an instructional guide and may be adjusted as the semester progresses. Any updates will be announced via Canvas.}

\end{document}
